\section{Irreducible residual inertia at the first wild index}\label{sec:irreducible}
The comparison with irreducible residual inertia can be settled by the ramification of one point at a time. This keeps the field degrees separate from the dimensions of congruence layers.

\begin{theorem}\label{thm:irreducible-recovery}
Let $p\ge5$, let $K$ have algebraically closed residue field, $e_K=p(p-1)$, and $\mu_p\subset K$. If two height-two groups with potentially supersingular good reduction have isomorphic $p^3$-torsion and their common residual inertia representation is irreducible over $\F_p$, then their finite special inertia characters agree.
\end{theorem}

Put $q=p^2$, $N=q-1$, and $e=p(p-1)$. Take the common tame good field $L/K$ of degree $N$. We argue by contradiction, assuming that the special characters differ.

\subsection{Leading coordinates and common point fields}
The residual wild image is a normal $p$-subgroup of the residual inertia image. Its nonzero fixed space is stable under inertia, so irreducibility makes that subgroup trivial. The residual representation is tame of degree two; its order divides $N$, and its whole first-level module is rational over $L$. A two-slope Newton polygon would give a canonical invariant high line. The polygon therefore has one slope. In the Honda coordinate its nonzero $p^j$-torsion valuations over $L$ are
\begin{equation}\label{eq:irreducible-values}
                              \frac e{q^{j-1}}.
\end{equation}
The first value is $v_L(p)/N=e$; every inverse polygon has a single segment and divides the value by $q$.

Let $c_i$ be the special tangent exponents. The first-level leading-coordinate spaces have tame exponents $c_i+e$. Each is a scalar multiple of $\F_q$: the leading-coordinate map is injective on the two-dimensional residual module, and its irreducible tame scalar generates $\F_q$ over $\F_p$. An equivariant isomorphism between two such spaces has, after scalar normalization, either the form $z\mapsto z$ or the form $z\mapsto z^p$. The scalar form gives equal special exponents. Under our contradiction hypothesis the form must therefore be
\begin{equation}\label{eq:frobenius-orientation}
                                  z_2=z_1^p.
\end{equation}
It remains true after changing the uniformizer of a point field, which only rescales the two leading spaces. At the boundary where the coefficient of $T^p$ could lie on the first Newton segment, the additive initial polynomial would contain a nonzero $T^p$ term. It cannot vanish on a scalar multiple of all of $\F_q$. Thus this boundary is excluded by irreducibility; the later leading multiplication terms are uniquely the degree-$q$ terms.

We use the following elementary ramification calculation. Let $T/S$ be a totally ramified Galois extension in mixed characteristic, with uniformizer $\varpi$. Put
\[
 b(g)=v_T(g\varpi-\varpi)-1.
\]
Expand a coordinate $x=\sum_{j\ge A}[a_j]\varpi^j$ using residue representatives fixed by $g$. Then $v_T(gx-x)\ge A+b(g)$, with equality when $p\nmid A$. Below $v_T(p)$, if $v_p(A)=1$, the first possible terms of the change have degrees
\begin{equation}\label{eq:first-change}
                         A+pb(g),\qquad J+b(g),
\end{equation}
where $J$ is the first nonzero exponent prime to $p$, possibly infinite. To verify this, write $g\varpi=\varpi(1+d\varpi^b+\cdots)$ and expand each monomial. For $j=p^tu$, $p\nmid u$, the first term below $v_T(p)$ is $u[a_j]d^{p^t}\varpi^{j+p^tb}$. Among the exponents divisible by $p$, the leading exponent $A$ is earliest; among the other exponents, $J$ is earliest. If the two displayed degrees differ, the smaller cannot cancel. If they agree, their leading coefficients must be included together. This proves every use of the calculation below.

Write
\[
 F_2=L(M_i[p^2]),\qquad F_3=L(M_i[p^3]).
\]
The torsion isomorphism makes these the same fields for $i=1,2$. Over $L$ the matrices are congruent to $1$ modulo $p$. Let $W_1,W_2\subset M_2(\F_p)$ be the first and terminal congruence layers, and write $r_j=\dim W_j$. Matrix multiplication gives
\begin{equation}\label{eq:matrix-powers}
 (1+pA+p^2B)^p=1+p^2A\pmod{p^3},\qquad W_1\subset W_2.
\end{equation}

\begin{lemma}\label{lem:full-point-orbits}
For every primitive $p^2$-point $P$ and every $p^3$-point $P^+$ with $pP^+=P$, the common point fields $E=L(P)$ and $T=L(P^+)$ satisfy
\[
                [E:L]=q,\qquad[T:E]=q,\qquad[T:L]=q^2.
\]
The two successive extensions are Galois, with the full $p$-torsion translation group. Moreover $W_1v=\F_p^2$ for every $v\ne0$.
\end{lemma}
\begin{proof}
Since $pP$ and every $p$-torsion point are rational over $L$, all conjugates of $P$ lie in $L(P)$. Thus $E/L$ is Galois, of degree at most $q$. The denominator of $v_L(P)=(p-1)/p$ makes its degree at least $p$. Similarly $T/E$ is Galois of degree at most $q$.

Suppose $[E:L]=p$. The denominator of $v_L(P^+)=(p-1)/p^3$ then forces $[T:L]=p^3$ and $[T:E]=q$. In $T$ the two coordinates $x_i$ have valuation $A=p-1$, prime to $p$, and every nonidentity translation displaces them by a nonzero $p$-torsion coordinate of valuation $D=ep^3$. The exact prime-valuation calculation gives one lower break $D-A$ and leading differences
\[
       \overline{(gx_i-x_i)/\varpi^D}=\overline A\,a_i d_g,
       \qquad a_i=\overline{x_i/\varpi^A}\ne0.
\]
These are below $v_T(p)=ND$. The two leading translation maps are scalar multiples of the same $d_g$, although their values run through the whole two-dimensional space and are related by~\eqref{eq:frobenius-orientation}. This would make Frobenius on $\F_q$ scalar, which it is not. Hence $[E:L]=q$. The orbit is $W_1v$, proving its asserted surjectivity. Equation~\eqref{eq:matrix-powers} then makes terminal matrices supply every $p$-torsion translation of $P^+$. Thus $[T:E]=q$ as well.
\end{proof}

\begin{lemma}\label{lem:three-first-dimensions}
One has $r_1\ge3$.
\end{lemma}
\begin{proof}
If $r_1=2$, evaluation at any nonzero vector is an isomorphism by Lemma~\ref{lem:full-point-orbits}. The stabilizer of $P^+$ is then contained in the terminal matrix subgroup, which is central modulo $p^3$. Thus $T/L$ is Galois. Put $G=\Gal(T/L)$ and $U=\Gal(T/E)$. Then
\[
               |G|=p^4,\qquad U\simeq\F_p^2,
                    \qquad G/U\simeq\F_p^2,
\]
and the $p$th-power map $G\to U$ is onto by~\eqref{eq:matrix-powers}.

In $T$ the three relevant coordinate valuations are
\[
 A=e,\qquad D_1=qA,\qquad D_2=q^2A,\qquad v_T(p)=ND_2.
\]
Elements outside $U$ displace the selected point by order-$p^2$ points, of valuation $D_1$; nonzero elements of $U$ displace it by order-$p$ points, of valuation $D_2$. Normalize each leading space using the selected points $pP^+$ and $p^2P^+$ as parameter $1$. If the primitive leading coefficient is $a_i$, the displacement coefficients have the forms
\[
            a_i^q\lambda_i(g),\qquad a_i^{q^2}\lambda_i(h),
                        \qquad\lambda_2=\lambda_1^p.
\]
The parameter maps on $G/U$ and $U$ are additive isomorphisms onto $\F_q$. The relation on the second-level leading spaces follows from the first-level relation by taking the unique $q$th roots.

Let $a$ be the smallest lower break of $G$. Every nonzero element of $U$ has larger break, since it is a $p$th power and the first additive ramification coefficient of a $p$th power is zero. Set
\[
 a_0=(D_1-A)/p=N(p-1),\qquad
 j_0=A+(p-1)a_0=D_1-a_0,\qquad j_0\equiv-1\pmod p.
\]
Consider an element outside $U$ of break $a$ in the two expansions~\eqref{eq:first-change}. If $a<a_0$, cancellation below $D_1$ forces the first ordinary degree to be $J_i=A+(p-1)a$, with nonzero coefficient. If $a>a_0$, achieving degree $D_1$ forces $J_i=D_1-a$, again with nonzero coefficient. An exponent divisible by $p$ in either assertion is already impossible. At every terminal break $b>a$, that fixed ordinary term is uniquely earliest: the differences from the divisible term are respectively $(p-1)(b-a)$ and $(p-1)(b-a)+p(a-a_0)$. Consequently $J_i+b=D_2$, and both terminal translation maps become scalar multiples of the common $d_h$. This contradicts their Frobenius relation. Therefore $a=a_0$.

No ordinary degree below $j_0$ can now occur. Write $f_i$ for its coefficient, allowing zero. At degree $D_1$, for a minimum-break element, the leading equality is
\begin{equation}\label{eq:boundary-leading}
 a_i^q\lambda_i(g)=(A/p)a_i d_g^p+j_0f_i d_g.
\end{equation}
If one $f_i\ne0$, every terminal break equals $b_*=D_2-j_0>a_0$, and this ordinary term is uniquely first. In the other model the same break also forces a nonzero term at exactly $j_0$, because
\[
 b_*-(D_2-A)/p=(p-1)(D_2-D_1)/p>0.
\]
We again obtain scalar proportionality on $U$, a contradiction. Thus both $f_i$ vanish. Every nonterminal element must have break $a_0$, since a larger break would put every coordinate change after $D_1$. The first ramification coefficient identifies $G/U$ with a two-dimensional additive space, and~\eqref{eq:boundary-leading} becomes
\[
                       \lambda_i(g)=(A/p)a_i^{1-q}d_g^p.
\]
The two maps are scalar proportional on all of $\F_q$, again contradicting their Frobenius relation. Hence $r_1=2$ is impossible.
\end{proof}

\subsection{The two point conductors}
\begin{lemma}\label{lem:point-conductors}
The last upper breaks of $E/L$ and $T/E$ are respectively
\[
             c_1=N(p-1),\qquad c_2=(q^2-1)(p-1).
\]
The last upper break of $F_2/L$ is $c_1$, and that of $F_3/F_2$ is at most $p^{r_1-2}c_2$.
\end{lemma}
\begin{proof}
In either successive point extension the group is $\F_p^2$, the coordinate valuation is $A=e$, and every nonidentity element has displacement $D$, with $D=qA$ or $q^2A$. The two leading translation maps fill $\F_q$ and have the Frobenius relation. Put $b_0=(D-A)/p$.

If there is just one lower break $b<b_0$, cancellation in~\eqref{eq:first-change} forces $J=A+(p-1)b$ and $d_g^{p-1}=\kappa\ne0$ on all $q-1$ nonzero leading ramification coefficients. This is impossible for a polynomial of degree $p-1$. If $b>b_0$, the unique ordinary term of degree $D-b$ makes the leading maps scalar proportional, also impossible. Hence a single break equals $b_0$.

If there are two lower breaks $t<u$, then $t<b_0$. For if $t\ge b_0$, no ordinary term can precede degree $D$ at $t$, and every term would follow $D$ at $u$. Cancellation at $t$ therefore forces $J=A+(p-1)t$. At $u$ it is uniquely earliest, so
\[
                               u+(p-1)t=D-A.
\]
Herbrand's formula makes the last upper break $(u+(p-1)t)/p=b_0$. This proves the two stated values.

Choose a basis $P_1,P_2$ of $p^2$-torsion. Then $F_2=L(P_1)L(P_2)$. Both elementary abelian point fields have last break $c_1$, so their compositum has that last break: every cyclic character is a sum of characters coming from the factors, and conductors do not increase under addition.

Lift the basis to $P_i^+$ and set $E_i=L(P_i)$, $T_i=L(P_i^+)$ and $D_i=F_2T_i$. The terminal translations are already the full translations of $T_i/E_i$, giving
\[
 F_2\cap T_i=E_i,\quad[D_i:F_2]=q,\quad
 [F_2:E_i]=p^{r_1-2},\quad D_1D_2=F_3.
\]
For completeness let $H_i$ and $S_i$ be the stabilizers of $P_i$ and $P_i^+$ in $\Gal(F_3/L)$, and let $U$ be the terminal subgroup. Surjectivity of terminal translations says $US_i=H_i$, which proves the intersection. The subgroup $U\cap S_i$ is central. For first-layer matrices $A,C$ of $H_i$, both kill the residual vector $v_i$, so $[A,C]v_i=0$; their $p$th powers also kill it. Hence $D_i/E_i$ is an elementary abelian Galois extension.

For upper index $u>c_2$, its ramification group has trivial image in $\Gal(T_i/E_i)$. Its intersection with $\Gal(D_i/F_2)$ is therefore trivial, since $D_i=F_2T_i$. The subgroup form of Herbrand's theorem bounds the last break of $D_i/F_2$ by $\psi_{F_2/E_i}(c_2)\le p^{r_1-2}c_2$. The two fields $D_i$ generate the elementary abelian terminal extension, proving the same bound for $F_3/F_2$. These are the quotient and subgroup conversions of \citet[Chapter~IV, \S3, Propositions~14--15, pp.~75--76]{Serre1979}.
\end{proof}

\subsection{Contradiction from the determinant}
We finish the proof of Theorem~\ref{thm:irreducible-recovery}. On the good field $L$, the determinant is cyclotomic. First suppose that the trace on $W_1$ is zero. Lemma~\ref{lem:three-first-dimensions} gives $W_1=\mathfrak{sl}_2$ and $r_1=3$. Its determinant is trivial modulo $p^2$, so $\mu_{p^2}\subset L$.

Let $K_0$ be the Cohen field of $K$. The extension $K_0(\mu_{p^2})/K_0(\mu_p)$ has degree $p$. Since $K/K_0(\mu_p)$ also has degree $p$ and $L/K$ is tame, adjoining $\mu_{p^2}$ to $K$ cannot have degree $p$. Thus $K=K_0(\mu_{p^2})$. The cyclotomic extension $C=L(\mu_{p^3})$ has degree $p$ over $L$ and break $N^2$: before the tame base change its break is $p^2-1=N$, and the tame degree is $N$.

Because $N^2>c_1$, $C\cap F_2=L$. In the elementary abelian compositum its cyclotomic character has the unique largest break. After restriction to $F_2$ its conductor is exactly
\[
 \psi_{F_2/L}(N^2)
       =\psi_{F_2/L}(c_1)+p^3(N^2-c_1)>p^3(N^2-c_1).
\]
It is a nonzero terminal character, since $C\subset F_3$. Yet Lemma~\ref{lem:point-conductors} bounds it by $pc_2$, contradicting
\[
 p^3(N^2-c_1)-pc_2=pN(p-1)(p^3-p^2-1)>0.
\]
The trace on $W_1$ is consequently nonzero.

It follows that $\mu_{p^2}\not\subset K$: otherwise the first congruence determinant over $L$ would be trivial. Apply Theorem~\ref{thm:cyclotomic-exact} to $K(\mu_{p^3})/K$, with last upper break $\delta_2$. It gives
\[
 d_2=N\delta_2\ge Np(p-1)>c_1.
\]
The intersection of the cyclic degree-$p^2$ field $L(\mu_{p^3})$ with the elementary abelian field $F_2$ is exactly $L(\mu_{p^2})$. Its nontrivial terminal character consequently has conductor
\[
 \psi_{F_2/L}(d_2)
     =\psi_{F_2/L}(c_1)+p^{r_1}(d_2-c_1)
     >p^{r_1}(d_2-c_1).
\]
One may verify the equality directly from upper ramification groups: above $c_1$ the quotient $F_2/L$ is trivial, while the cyclotomic subgroup remains nontrivial up to $d_2$; the subgroup conversion is therefore exact. But
\[
 q(d_2-c_1)-c_2\ge N(p-1)(p^3-2p^2-1)>0.
\]
This contradicts the terminal bound $p^{r_1-2}c_2$ once more. The assumption of different characters was impossible, proving the theorem.\hfill$\square$
